% Moved 2026-10-28 from dl_rnn.tex: the closing summary, to be input last.
%%%%%%%%%%%%%%%%%%%%%%%%%%%%%%%%%%%%%%%%%%%%%%%%%%%
\begin{frame}[fragile] \frametitle{Summary}
\begin{itemize}
% [upgrade-deck: rewritten 2026-10-26 | Tasks 1 and 5c | grammar, "solve" overclaims (see Fading by Numbers), long bullet split; original below]
% \item RNNs suffer from vanishing gradients and caused by long series of multiplications of small values, diminishing the gradients and causing the learning process to become degenerate.
% \item LSTMs solve the problem using a unique additive gradient structure that includes direct access to the forget gate’s activations, enabling the network to encourage desired behaviour from the error gradient using frequent gates update on every time step of the learning process.
% [Rewritten 2026-10-28 | author: the deck now has separate RNN, GRU and LSTM sections | the summary covers all three; original above]
\item \textbf{RNN:} one state, the same weights at every step, so any length and order kept
  \begin{itemize}
  \item Vanishing gradients: long series of multiplications of small values, so the early steps barely learn
  \end{itemize}
\item \textbf{GRU:} update and reset gates choose between keeping and rewriting the memory
  \begin{itemize}
  \item A short path for the old memory; still one vector for memory and output
  \end{itemize}
\item \textbf{LSTM:} a separate cell state with forget, input and output gates
  \begin{itemize}
  \item The additive cell-state path gives the gradient a route that does not shrink when the forget gate stays near 1
  \end{itemize}
\item Gates help, but they do not remove the problem: at the untrained start a GRU or an LSTM also fades
\item To get one vector from a sequence: last state, pooling over time, attention, or both directions
\end{itemize}

{\tiny (Ref: How LSTM networks solve the problem of vanishing gradients
- Nir Arbel)}

\end{frame}
